\chapter{Microlocal Spectrum Condition And Hadamard Geometry}
\label{ch:curved-microlocal-spectrum-condition}

\ConventionsMixed

This chapter develops the microlocal form of the positive-frequency condition
on curved spacetimes.  The replacement for translation-invariant spectral
support is a wavefront-set condition on distributions, and the local
geometric form of the two-point function is the Hadamard parametrix.

\section{Wavefront Set}

Let \(u\in\mathcal D'(X)\) be a distribution on a smooth manifold \(X\).  A
point \((x,k)\in T^*X\setminus\{0\}\) is absent from the wavefront set
\(\operatorname{WF}(u)\) if there is a cutoff \(\chi\in C_c^\infty(X)\) with
\(\chi(x)\ne0\) and a conic neighborhood \(\Gamma\) of \(k\) such that the
Fourier transform of \(\chi u\) decays faster than every power in \(\Gamma\).
The wavefront set records both the singular support and the covector
directions of oscillation.

\begin{definition}[Wavefront set in a chart]
\label{def:wavefront-set-chart}
Let \(u\in\mathcal D'(X)\).  In a coordinate chart \(U\subset X\), write
\(\widehat{\chi u}(\xi)\) for the Fourier transform of the compactly
supported distribution \(\chi u\), \(\chi\in C_c^\infty(U)\), in the chosen
coordinates.  A nonzero covector \(k\in T_x^*X\) is a regular direction for
\(u\) if there exist \(\chi\in C_c^\infty(U)\), \(\chi(x)\ne0\), and an open
conic neighborhood \(\Gamma\subset\mathbb R^D\setminus\{0\}\) of the
coordinate representative of \(k\) such that, for every \(N\),
\[
  |\widehat{\chi u}(\xi)|
  \leq C_N(1+|\xi|)^{-N},
  \qquad
  \xi\in\Gamma .
  \label{eq:microlocal-rapid-decay-direction}
\]
The wavefront set \(\operatorname{WF}(u)\) is the complement, in
\(T^*X\setminus0\), of the regular directions.
\end{definition}

The definition is independent of the chart.  Under a change of coordinates,
the localized Fourier transform is an oscillatory integral whose phase has
non-degenerate mixed Hessian at the point under consideration; integration
by parts in directions where the transformed covector is separated from the
old singular cone preserves rapid decay.  Thus \(\operatorname{WF}(u)\) is a
closed conic subset of \(T^*X\setminus0\).

\begin{theorem}[Product criterion for distributions]
\label{thm:microlocal-product-criterion}
Let \(u,v\in\mathcal D'(X)\).  Suppose that there is no
\[
  (x,k)\in\operatorname{WF}(u)
  \quad\hbox{and}\quad
  (x,-k)\in\operatorname{WF}(v).
  \label{eq:microlocal-no-opposite-covectors}
\]
Then the pointwise product \(uv\) has a unique distributional extension
which agrees with the ordinary product when one factor is smooth.  Moreover
\[
\begin{aligned}
  \operatorname{WF}(uv)
  \subset&
  \operatorname{WF}(u)\cup\operatorname{WF}(v)\\
  &\cup
  \{(x,k+\ell):(x,k)\in\operatorname{WF}(u),\
        (x,\ell)\in\operatorname{WF}(v),\ k+\ell\ne0\}.
\end{aligned}
  \label{eq:microlocal-product-wavefront-bound}
\]
\end{theorem}

\begin{proof}
The statement is local, so work in one coordinate chart and multiply by a
cutoff equal to \(1\) near the point.  If \(u\) and \(v\) are compactly
supported distributions, the formal Fourier transform of the product is the
convolution
\[
  \widehat{uv}(\xi)
  =
  (2\pi)^{-D}\int \widehat u(\eta)\widehat v(\xi-\eta)\,\dd\eta .
  \label{eq:microlocal-product-convolution}
\]
The obstruction to defining this integral distributionally at large
\(|\eta|\) is simultaneous non-decay of \(\widehat u(\eta)\) and
\(\widehat v(\xi-\eta)\).  As \(|\eta|\to\infty\), the second direction is
asymptotic to \(-\eta\).  Thus a large-frequency obstruction can occur only
from an opposite pair of singular covectors \((k,-k)\), which is excluded by
\eqref{eq:microlocal-no-opposite-covectors}.

Choose a conic partition of unity in the \(\eta\)-variable.  On pieces where
\(\eta\) avoids \(\operatorname{WF}(u)\), the factor \(\widehat u(\eta)\) is
rapidly decreasing after localization.  On pieces where \(\xi-\eta\) avoids
\(\operatorname{WF}(v)\), the factor \(\widehat v(\xi-\eta)\) is rapidly
decreasing.  The exclusion of opposite cones gives such a finite conic cover
of the large-\(\eta\) region.  The bounded-\(\eta\) region is harmless
because Fourier transforms of compactly supported distributions have at most
polynomial growth.  This defines
\eqref{eq:microlocal-product-convolution}, and the same conic decomposition
proves rapid decay of \(\widehat{uv}\) away from the three sets displayed in
\eqref{eq:microlocal-product-wavefront-bound}.

Uniqueness follows by regularizing \(u\) and \(v\) with a smoothing
approximate identity.  The estimates above show convergence in every allowed
conic region and independence of the regularization.  If one factor is
smooth, the construction reduces to ordinary multiplication because the
smooth factor has rapidly decreasing localized Fourier transform.
\end{proof}

This criterion is the local analytic reason that Wick products and
time-ordered products require subtraction before coincident points are taken.

\section{Hadamard Two-Point Condition}

Let \(M\) be a globally hyperbolic Lorentzian spacetime and let \(\phi\) be a
real scalar field satisfying
\[
  P\phi=0,\qquad P=\Box_g+m^2+\xi R .
\]
A quasifree state is determined by its two-point function
\[
  \omega_2(x,y)=\omega(\phi(x)\phi(y)).
\]
The time orientation on tangent vectors induces a convention on covectors:
a nonzero causal covector \(k\in T_x^*M\) is future directed when its metric
dual
\[
  k^\sharp=g^{\mu\nu}k_\nu\partial_\mu
\]
is a future-directed causal vector.  In mostly-plus Minkowski coordinates,
this means that a future null covector has \(k_0<0\).  This convention makes
the positive-frequency plane wave \(e^{-\ii Et+\ii{\bf p}\cdot{\bf x}}\)
carry future null covector \((-E,{\bf p})\).

The microlocal Hadamard condition is
\[
  \operatorname{WF}(\omega_2)
  =
  \{(x,k_x;y,-k_y): (x,k_x)\sim(y,k_y),\ k_x\in\overline V_+^*\}.
\]
Here \((x,k_x)\sim(y,k_y)\) means that \(x\) and \(y\) are joined by a null
geodesic and that \(k_y\) is the parallel transport of \(k_x\) along that
geodesic.  The cone \(\overline V_+^*\) is the closed future causal cone in
the cotangent space.

The minus sign on the second covector is not a convention that can be
dropped.  It records that a two-point distribution is a distribution on
\(M\times M\).  If the same future covector is transported from \(x\) to
\(y\), the corresponding singular phase is locally of the form
\[
  \exp\bigl(\ii\lambda(\Phi(x)-\Phi(y))\bigr),
  \qquad \lambda\to+\infty,
\]
so the covector on the second copy of \(M\) is the negative of the
transported covector.

\section{Higher Microlocal Spectrum Condition}

For non-Gaussian states, the two-point condition is not enough.  The
curved-spacetime replacement of the Wightman spectral condition is a family
of cone constraints on all \(n\)-point distributions.

\begin{definition}[Microlocal spectrum cone]
\label{def:microlocal-spectrum-cone}
Let \(x_1,\ldots,x_n\in M\).  A covector tuple
\[
  (x_1,k_1;\ldots;x_n,k_n)\in T^*M^n\setminus0
\]
belongs to \(\Gamma_n(M)\) if there is a finite directed graph with vertex
set \(\{1,\ldots,n\}\) such that each edge \(e:i\to j\) is assigned a
future-directed causal covector field \(p_e\) parallel transported along a
future-directed causal curve from \(x_i\) to \(x_j\), and
\[
  k_i
  =
  \sum_{e:i\to j}p_e(x_i)
  -
  \sum_{e:j\to i}p_e(x_i)
  \qquad (i=1,\ldots,n).
  \label{eq:microlocal-spectrum-graph-balance}
\]
The empty graph gives the zero tuple, which is omitted from
\(T^*M^n\setminus0\).
\end{definition}

\begin{definition}[Microlocal spectrum condition for a state]
\label{def:microlocal-spectrum-condition-state}
A state \(\omega\) on a scalar field algebra satisfies the microlocal
spectrum condition when all \(n\)-point distributions
\[
  \omega_n(f_1,\ldots,f_n)
  =
  \omega(\phi(f_1)\cdots\phi(f_n))
\]
obey
\[
  \operatorname{WF}(\omega_n)\subset\Gamma_n(M).
  \label{eq:microlocal-spectrum-condition-all-n}
\]
\end{definition}

For \(n=2\), a graph with one edge \(1\to2\) gives
\[
  (k_1,k_2)=(p,-p),
\]
which is the sign pattern in the Hadamard two-point condition.  For
quasifree Hadamard states, Wick's theorem reduces every \(n\)-point
distribution to a finite sum of products of two-point functions.  Applying
Theorem~\ref{thm:microlocal-product-criterion} and the graph description
above shows that the microlocal spectrum condition for all \(n\) follows
from the two-point Hadamard condition.

\section{Hadamard Parametrix}

In a geodesically convex neighborhood, the singular part of a four-dimensional
Hadamard two-point function has the form
\[
  H_\epsilon(x,y)
  =
  \frac1{8\pi^2}
  \left[
    \frac{U(x,y)}{\sigma_\epsilon(x,y)}
    +
    V(x,y)\log\!\left(\frac{\sigma_\epsilon(x,y)}{\ell^2}\right)
  \right],
\]
where \(\sigma(x,y)\) is one half the squared geodesic distance,
\[
  \sigma_\epsilon(x,y)=\sigma(x,y)+\ii\epsilon(t(x)-t(y))+\epsilon^2,
\]
\(U=\Delta^{1/2}\) is the square root of the van Vleck determinant, \(V\) is
determined recursively by the wave operator \(P\), and \(\ell\) is a length
scale.  A Hadamard state has
\[
  \omega_2(x,y)-H_\epsilon(x,y)
\]
smooth in the neighborhood, after the distributional boundary value is taken.

\section{Propagation Of Singularities}

For a real principal type operator \(P\), singularities of solutions to
\(Pu=0\) propagate along the Hamiltonian flow of the principal symbol.  For
the Klein--Gordon operator the principal symbol is
\[
  p(x,k)=g^{\mu\nu}(x)k_\mu k_\nu .
\]
The Hamiltonian vector field on \(T^*M\) is
\[
  H_p
  =
  \frac{\partial p}{\partial k_\mu}\frac{\partial}{\partial x^\mu}
  -
  \frac{\partial p}{\partial x^\mu}\frac{\partial}{\partial k_\mu}.
\]
Thus its integral curves satisfy
\[
  \dot x^\mu=2g^{\mu\nu}(x)k_\nu,
  \qquad
  \dot k_\mu=-(\partial_\mu g^{\alpha\beta})(x)k_\alpha k_\beta .
  \label{eq:microlocal-kg-hamilton-equations}
\]
Let \(v^\mu=\dot x^\mu\).  Then
\[
  v^\mu=2k^{\sharp\mu},
  \qquad
  g_{\mu\nu}v^\mu v^\nu=4p(x,k).
\]
On \(p=0\), the projected curve \(x(s)\) is therefore null.  Lowering the
index in the first Hamilton equation gives \(k_\mu=\frac12g_{\mu\nu}v^\nu\).
Differentiate this identity and use the second Hamilton equation together
with
\[
  \partial_\mu g^{\alpha\beta}
  =
  -g^{\alpha\rho}g^{\beta\sigma}\partial_\mu g_{\rho\sigma}.
\]
One obtains
\[
  g_{\mu\nu}\left(\ddot x^\nu+\Gamma^\nu_{\rho\sigma}
       \dot x^\rho\dot x^\sigma\right)=0 .
\]
Since \(g\) is non-degenerate, \(x(s)\) is an affinely parametrized null
geodesic, and \(k\) is parallel transported along it up to the fixed factor
\(1/2\).

Hence once the two-point function has positive-frequency wavefront
directions on one Cauchy surface, the equation of motion propagates those
directions along null geodesics.  This is the geometric content of the
microlocal spectrum condition.

\section{Wick Products}

For a Hadamard state, define the Wick square by point splitting:
\[
  :\!\phi^2\!:_H(f)
  =
  \int_M f(x)
  \lim_{y\to x}
  \bigl(\phi(x)\phi(y)-H(x,y)\mathbf 1\bigr)\,\dd\operatorname{vol}_g(x).
\]
Changing the parametrix scale or adding a smooth local bisolution changes the
Wick square by a local curvature polynomial:
\[
  :\!\phi^2\!:'_H
  =
  :\!\phi^2\!:_H
  +
  c_1R+c_2m^2 .
\]
Thus the local covariant field is defined only after the finite local
renormalization coordinates are specified.

\section{Time-Ordered Products}

The time-ordered two-point distribution has a wavefront set with both
future and past directions arranged according to time ordering.  Products of
time-ordered distributions at coincident points violate the wavefront-set
product criterion.  Perturbative algebraic renormalization extends these
distributions from configuration space with diagonals removed to the full
configuration space.  The extension freedom is local and covariant, with
degree bounded by scaling.  Contact terms in curved spacetime are precisely
the local distributions supported on these diagonals, with covariance and
Ward identities imposed as constraints.

\begin{openproblem}[Interacting microlocal spectrum condition]
\label{op:interacting-microlocal-spectrum-condition}
Construct interacting locally covariant QFTs on curved backgrounds with
Hadamard states, local Wick and time-ordered products, stress tensor,
background gauge fields, and anomaly lines, and prove the microlocal spectrum
condition for the resulting nonperturbative Wightman-type distributions.
\end{openproblem}
